% Adaptación de FORMATO.tex para ACD2; original sin cambios.
\documentclass[12pt,a4paper]{report}

% ------------------------------------------
%   PAQUETES BÁSICOS
% ------------------------------------------
\usepackage[spanish]{babel}
\usepackage[utf8]{inputenc}
\usepackage[T1]{fontenc}
\usepackage{graphicx}
\usepackage{amsmath, amssymb}
\usepackage{booktabs}
\usepackage{geometry}
  \geometry{left=3cm, right=2.5cm, top=3cm, bottom=2.5cm}
\usepackage{setspace}
  \onehalfspacing
\usepackage{titlesec} 
  \titleformat{\chapter}[block]{\bfseries\large}{\thechapter.}{1em}{}
  \titleformat{\section}[block]{\bfseries}{\thesection.}{1em}{}
\usepackage{tocloft}
  \setlength{\cftbeforechapskip}{6pt}
  \renewcommand{\cftchapfont}{\bfseries}
\usepackage{longtable}
\usepackage{array}
\usepackage{verbatim}

\usepackage{cite}
\usepackage{hyperref}
% Soporte mínimo para tablas, rutas y ejemplos HTTP.
\usepackage{calc}
\usepackage{xurl}
\usepackage{fvextra}
\providecommand{\tightlist}{\setlength{\itemsep}{0pt}\setlength{\parskip}{0pt}}
\setlength{\emergencystretch}{3em}

\addto\captionsspanish{
    \renewcommand{\contentsname}{Índice General}
    \renewcommand{\tablename}{Tabla}
}

% ------------------------------------------
%   FONDO DE PÁGINA (SOLO CARÁTULA)
% ------------------------------------------
\usepackage{eso-pic}
\newcommand\BackgroundPic{%
  \put(0,0){%
    \parbox[b][\paperheight]{\paperwidth}{%
      \vfill
      \centering
      \IfFileExists{fondo.jpeg}{\includegraphics[width=\paperwidth,height=\paperheight]{fondo.jpeg}}{\parbox{0.8\paperwidth}{\centering\small [PENDIENTE: fondo.jpeg original]}}%
      \vfill
}}}

% ------------------------------------------
%   ENCABEZADO CON LOGO (OTRAS PÁGINAS)
% ------------------------------------------
\usepackage{fancyhdr}
\pagestyle{fancy}
\fancyhf{}
\setlength{\headheight}{40pt}
\fancyhead[L]{\IfFileExists{logo_unl.jpeg}{\includegraphics[height=1.2cm]{logo_unl.jpeg}}{\fbox{\scriptsize PENDIENTE: logo\_unl.jpeg}}}
\fancyfoot[C]{\thepage}
\renewcommand{\headrulewidth}{0pt}

\fancypagestyle{plain}{
    \fancyhf{}
    \fancyhead[L]{\IfFileExists{logo_unl.jpeg}{\includegraphics[height=1.2cm]{logo_unl.jpeg}}{\fbox{\scriptsize PENDIENTE: logo\_unl.jpeg}}}
    \fancyfoot[C]{\thepage}
    \renewcommand{\headrulewidth}{0pt}
}

% ------------------------------------------
%   METADATOS DEL DOCUMENTO
% ------------------------------------------
\newcommand{\TituloTrabajo}{ACD2 -- Semana 2: Requisitos verificables y HTTP}
\newcommand{\Autor}{[PENDIENTE: integrantes]}
\newcommand{\Director}{[PENDIENTE: docente]}
\newcommand{\Carrera}{Carrera de Computación}
\newcommand{\Facultad}{Facultad de la Energía, las Industrias y los Recursos Naturales No Renovables}
\newcommand{\Universidad}{Universidad Nacional de Loja}
\newcommand{\Ciudad}{Loja, Ecuador}
\newcommand{\Fecha}{[PENDIENTE: fecha de la sesión]}

\begin{document}

\AddToShipoutPictureBG*{\BackgroundPic}

% ------------------------------------------
%   PORTADA
% ------------------------------------------
\begin{titlepage}
  \centering
  \vspace*{2cm}
  {\bfseries\Large \Universidad}\\[0.5cm]
  {\bfseries\large \Facultad}\\[0.5cm]
  {\bfseries\large \Carrera}\\[2cm]
  {\bfseries\Large \TituloTrabajo}\\[2.5cm]

  \begin{flushright}
    \setlength{\fboxsep}{8pt}
    \fbox{%
      \begin{minipage}{0.4\textwidth}\small
        Desarrollo de Software Distribuido.
      \end{minipage}%
    }
  \end{flushright}
  \vspace{1cm}

  {\bfseries \large AUTOR:}\\[0.5em]
  {\large \Autor}\\[2em]
  {\bfseries \large TUTOR:}\\[0.5em]
  {\large \Director}\\[2em]
  {\bfseries \large \Ciudad}\\[1em]
  {\bfseries \large \Fecha}
\end{titlepage}

% ------------------------------------------
%   DECLARACIÓN DE USO DE IA GENERATIVA
% ------------------------------------------
\newpage
\chapter*{Declaración de Delegación a IA Generativa}
\addcontentsline{toc}{chapter}{Declaración de uso de IA Generativa}
Se utilizó ChatGPT/Codex para inspección estática del repositorio, contraste de requisitos con código, consulta de documentación técnica y redacción de esta entrega. No se utilizaron resultados de ejecución nuevos ni se atribuyen al docente o a estudiantes observaciones no aportadas.

El equipo deberá revisar el contenido, completar los datos y asumir responsabilidad por la versión presentada. La revisión humana final, participación y retroalimentación de clase permanecen pendientes. No se declara supervisión humana ya completada ni se trasladan las declaraciones del informe anterior contenido en la plantilla.

\noindent\textbf{Declaración enviada por: [PENDIENTE: responsable real].}
\newpage
\tableofcontents
\clearpage
\chapter*{Ficha de la actividad}
\addcontentsline{toc}{chapter}{Ficha de la actividad}
\textbf{Modalidad:} Actividad grupal con evidencia individual.\par
\textbf{Proyecto:} SIGEAC / SEC-UNL.\par
\textbf{Ciclo y paralelo:} [PENDIENTE].\par
\textbf{Base de revisión:} commit \texttt{faf868ea}.\par
\textbf{Propósito:} Formular requisitos funcionales y no funcionales verificables a partir del problema del proyecto y comprender los fundamentos de HTTP como protocolo base para la comunicación entre clientes y servicios web.

Esta entrega es una actualización académica complementaria de la especificación inicial. Conserva sus identificadores y diferencia requisitos existentes, mejoras propuestas y decisiones pendientes. No sustituye la aprobación institucional. Los ejemplos HTTP son reconstrucciones sintéticas del código y no capturas de tráfico. No se arrancaron servicios ni se ejecutaron suites para elaborar el informe.
\chapter{Metodología de análisis de requisitos y
HTTP}\label{metodologuxeda-de-anuxe1lisis-de-requisitos-y-http}

\section{Fuentes y procedimiento}\label{fuentes-y-procedimiento}

Se partió del diagnóstico de \path|AUDITORIA_DOCUMENTACION_SDD.md| y se
contrastaron las fuentes del proyecto con controladores y DTO reales. El
procedimiento fue: revisar contexto y alcance; seleccionar requisitos
principales; hacer explícitas condiciones y resultados; relacionar
escenarios y aceptación; interpretar solicitudes/respuestas; registrar
correcciones y pendientes de socialización.

Fuentes de producto, relativas a la raíz del repositorio:

{\def\LTcaptype{none} % do not increment counter
\begin{longtable}[]{@{}
  >{\raggedright\arraybackslash}p{(\linewidth - 4\tabcolsep) * \real{0.3333}}
  >{\raggedright\arraybackslash}p{(\linewidth - 4\tabcolsep) * \real{0.3333}}
  >{\raggedright\arraybackslash}p{(\linewidth - 4\tabcolsep) * \real{0.3333}}@{}}
\toprule\noalign{}
\begin{minipage}[b]{\linewidth}\raggedright
Código
\end{minipage} & \begin{minipage}[b]{\linewidth}\raggedright
Archivo
\end{minipage} & \begin{minipage}[b]{\linewidth}\raggedright
Uso en la actividad
\end{minipage} \\
\midrule\noalign{}
\endhead
\bottomrule\noalign{}
\endlastfoot
D1 & \path|docs/requirements/spec.md| & Actores, alcance, FR-001--040,
RNF-001--019, criterios SC y CL-001--010. \\
D2 & \path|docs/requirements/domain-rules.md| & Estados, fórmulas,
responsabilidades y CL-011--017. \\
D3 & \path|docs/requirements/process-flows.md| & Procesos P01--P16 y
escenarios alternativos. \\
D4 & \path|specs/005-complete-evaluation-workflows/spec.md| &
FR-041--070 e historias de flujos completos. \\
D5 & \path|docs/testing/process-acceptance.md| & Criterios AC-P y
AC-X. \\
D6 & \path|docs/testing/traceability.md| & Relaciones entre requisitos,
historias, tareas y pruebas. \\
D7 & \path|docs/technical/contracts/openapi.yaml| & Contrato HTTP; debe
contrastarse con los controladores. \\
D8 & \path|docs/architecture/complete-workflows.md| & Contexto,
contenedores y límites del backend modular. \\
\end{longtable}
}

Fuentes ejecutables específicas:

{\def\LTcaptype{none} % do not increment counter
\begin{longtable}[]{@{}
  >{\raggedright\arraybackslash}p{(\linewidth - 2\tabcolsep) * \real{0.5000}}
  >{\raggedright\arraybackslash}p{(\linewidth - 2\tabcolsep) * \real{0.5000}}@{}}
\toprule\noalign{}
\begin{minipage}[b]{\linewidth}\raggedright
Código
\end{minipage} & \begin{minipage}[b]{\linewidth}\raggedright
Archivo exacto
\end{minipage} \\
\midrule\noalign{}
\endhead
\bottomrule\noalign{}
\endlastfoot
C1 &
\path|backend/src/main/java/ec/edu/unl/sec/selfassessment/api/CareerController.java| \\
C2 &
\path|backend/src/main/java/ec/edu/unl/sec/selfassessment/application/CareerAdministrationService.java| \\
C3 &
\path|backend/src/main/java/ec/edu/unl/sec/selfassessment/domain/Career.java| \\
C4 &
\path|backend/src/main/java/ec/edu/unl/sec/identityaccess/api/SessionController.java| \\
C5 &
\path|backend/src/main/java/ec/edu/unl/sec/identityaccess/infrastructure/SessionAuthenticationFilter.java| \\
C6 &
\path|backend/src/main/java/ec/edu/unl/sec/identityaccess/infrastructure/SessionCookie.java| \\
C7 &
\path|backend/src/main/java/ec/edu/unl/sec/evaluationmodel/api/EvaluationModelController.java| \\
C8 &
\path|backend/src/main/java/ec/edu/unl/sec/evidence/api/EvidenceSourceController.java| \\
C9 &
\path|backend/src/main/java/ec/edu/unl/sec/reporting/domain/DocumentaryProgressCalculator.java| \\
C10 &
\path|backend/src/main/java/ec/edu/unl/sec/shared/api/ApiProblemHandler.java| \\
C11 &
\path|backend/src/main/java/ec/edu/unl/sec/shared/api/EntityTagSupport.java| \\
C12 & \path|frontend/src/app/core/http/api-client.ts| \\
C13 & \path|frontend/proxy.conf.json| y \path|frontend/angular.json| \\
C14 &
\path|backend/src/main/java/ec/edu/unl/sec/identityaccess/infrastructure/SecurityConfiguration.java| \\
C15 &
\path|backend/src/main/java/ec/edu/unl/sec/identityaccess/infrastructure/JpaIdentityAccessRepository.java| \\
C16 &
\path|backend/src/main/java/ec/edu/unl/sec/selfassessment/api/CareerCatalogProblemHandler.java| \\
\end{longtable}
}

Los ejemplos HTTP se reconstruyen a partir de estas fuentes con valores
sintéticos. No son capturas de tráfico ni resultados de pruebas
realizadas para esta entrega. La revisión es documental y estática; no
se arrancaron servicios ni se modificó el código de negocio.

\section{Problema del proyecto y
necesidad}\label{problema-del-proyecto-y-necesidad}

El proyecto aborda la organización y trazabilidad de la autoevaluación
interna de carreras: configurar modelos, planificar responsables,
aportar información, revisar documentación y conservar las decisiones y
resultados. La especificación exige distinguir tres actos: recopilar
información, aceptar documentalmente una evidencia y valorar
académicamente un indicador.

La necesidad funcional documentada consiste en disponer de información
identificable por carrera, proceso, indicador, fuente y versión, con
acceso autorizado e historial. Sin esa separación, una carga o
aceptación documental podría interpretarse indebidamente como
cumplimiento académico; el proyecto establece explícitamente que tal
equivalencia no debe producirse.

\textbf{Límite de la evidencia:} no se encontraron entrevistas, cifras
de tiempos actuales ni un diagnóstico institucional que pruebe cómo se
desarrolla hoy el proceso fuera del sistema. La descripción anterior
resume el problema abordado por la especificación; no afirma que una
institución utilice actualmente hojas de cálculo, pierda documentos o
tenga una tasa de errores determinada.

\section{Actores y responsabilidades}\label{actores-y-responsabilidades}

{\def\LTcaptype{none} % do not increment counter
\begin{longtable}[]{@{}
  >{\raggedright\arraybackslash}p{(\linewidth - 2\tabcolsep) * \real{0.5000}}
  >{\raggedright\arraybackslash}p{(\linewidth - 2\tabcolsep) * \real{0.5000}}@{}}
\toprule\noalign{}
\begin{minipage}[b]{\linewidth}\raggedright
Actor
\end{minipage} & \begin{minipage}[b]{\linewidth}\raggedright
Responsabilidad pertinente
\end{minipage} \\
\midrule\noalign{}
\endhead
\bottomrule\noalign{}
\endlastfoot
Administrador funcional & Mantener catálogos, modelos y asignaciones
dentro de sus capacidades autorizadas. \\
Coordinador de autoevaluación & Preparar y planificar procesos;
consultar avance y tramitar decisiones autorizadas. \\
Responsable documental & Crear borradores, aportar contenido y remitir
evidencias de su ámbito. \\
Revisor documental & Decidir sobre coberturas de una versión remitida,
conservando fundamento. \\
Evaluador académico & Valorar indicadores con reglas y sustentos
autorizados; no modificar revisión documental. \\
Consultor autorizado & Consultar resultados e historial dentro del
ámbito otorgado. \\
Administrador técnico & Operar configuración y cuentas según sus
capacidades, sin adquirir competencia académica. \\
Cuenta técnica & Ejecutar tareas técnicas explícitas; no adoptar
decisiones humanas por privilegio técnico. \\
\end{longtable}
}

Las autoridades, comisiones y funciones de aprobación de P12--P16 están
descritas como configurables. Su competencia real sigue pendiente de
validación institucional. No se equipara ``tener sesión'' con ``poder
aprobar''.

\section{Alcance y exclusiones}\label{alcance-y-exclusiones}

El MVP documental comprende P01--P11 y consulta histórica básica de P16:
acceso, carreras, modelos, autorización del proceso, planificación,
evidencias, revisión, subsanación y seguimiento documental. Valoración
académica, informes, mejora y cierre integral corresponden al alcance
posterior D7--D8, aunque el repositorio ya contiene implementación y
evidencias técnicas de esos módulos.

Se excluyen acreditación oficial, evaluación externa en el alcance
inicial, datos reales no autorizados, integraciones institucionales
productivas, aplicación móvil y operación sin conexión durante el MVP.
La disponibilidad multirregión y el despliegue productivo tampoco son
compromisos de esta actividad.

Para ACD2 se analizan escenarios representativos del MVP; la valoración
se incluye para demostrar la separación entre aceptación documental y
resultado académico.

\section{Clarificaciones de la semana
anterior}\label{clarificaciones-de-la-semana-anterior}

Se recuperan decisiones existentes; no se inventan respuestas previas.

{\def\LTcaptype{none} % do not increment counter
\begin{longtable}[]{@{}
  >{\raggedright\arraybackslash}p{(\linewidth - 6\tabcolsep) * \real{0.2500}}
  >{\raggedright\arraybackslash}p{(\linewidth - 6\tabcolsep) * \real{0.2500}}
  >{\raggedright\arraybackslash}p{(\linewidth - 6\tabcolsep) * \real{0.2500}}
  >{\raggedright\arraybackslash}p{(\linewidth - 6\tabcolsep) * \real{0.2500}}@{}}
\toprule\noalign{}
\begin{minipage}[b]{\linewidth}\raggedright
Registro
\end{minipage} & \begin{minipage}[b]{\linewidth}\raggedright
Pregunta recuperada o reformulada
\end{minipage} & \begin{minipage}[b]{\linewidth}\raggedright
Información documentada
\end{minipage} & \begin{minipage}[b]{\linewidth}\raggedright
Estado de confirmación
\end{minipage} \\
\midrule\noalign{}
\endhead
\bottomrule\noalign{}
\endlastfoot
CL-001 / CL-015 & ¿Qué función y autoridad pueden autorizar cada
actuación? & Funciones y política sintética; asignación y constancia
obligatorias. & Pendiente de validación institucional. \\
CL-002 & ¿Cuándo se permite combinar funciones o autorrevisar? &
Política de conflictos y excepciones temporales auditables. & Confirmar
autoridad y procedimiento de excepción. \\
CL-005 / CL-013 & ¿Qué escala o regla se aplica y qué puede
automatizarse? & Escalas configurables; cálculo solo con reglas
validadas. & Confirmar metodología de cada indicador. \\
CL-006 & ¿Cómo se calcula el avance? & Fuentes obligatorias aplicables;
métricas separadas. & Regla del piloto documentada; validación
institucional pendiente. \\
CL-008 & ¿Qué archivos y conservación se requieren? & Máximo inicial 25
MiB y conservación durante la vida del proceso. & Tipos, retención y
eliminación institucional por confirmar. \\
CL-011 / CL-014 & ¿Se utilizan cinco modelos completos y vigentes? &
Configuraciones representativas; procedencia y vigencia por corroborar.
& No confundir prueba sintética con catálogo normativo completo. \\
CL-016 / CL-017 & ¿Quién aprueba informes y qué permite cerrar? &
Políticas versionadas, pendientes y excepciones explícitas. & Sin
constancia institucional de aprobación. \\
\end{longtable}
}

No existe evidencia suficiente para afirmar que estas preguntas fueron
respondidas en la sesión anterior. El equipo debe completar persona,
fecha, respuesta y respaldo cuando disponga del registro real.

\chapter{Resultados de la actividad: especificación académica
revisada}\label{resultados-de-la-actividad-especificaciuxf3n-acaduxe9mica-revisada}

\section{Requisitos funcionales
principales}\label{requisitos-funcionales-principales}

Se conservan los IDs del proyecto. Las formulaciones siguientes precisan
la intención de sus fuentes para esta actividad; no introducen
funcionalidades ajenas ni cambian su aprobación institucional.

{\def\LTcaptype{none} % do not increment counter
\begin{longtable}[]{@{}
  >{\raggedright\arraybackslash}p{(\linewidth - 4\tabcolsep) * \real{0.3333}}
  >{\raggedright\arraybackslash}p{(\linewidth - 4\tabcolsep) * \real{0.3333}}
  >{\raggedright\arraybackslash}p{(\linewidth - 4\tabcolsep) * \real{0.3333}}@{}}
\toprule\noalign{}
\begin{minipage}[b]{\linewidth}\raggedright
ID y fuente
\end{minipage} & \begin{minipage}[b]{\linewidth}\raggedright
Formulación verificable
\end{minipage} & \begin{minipage}[b]{\linewidth}\raggedright
Actor y verificación
\end{minipage} \\
\midrule\noalign{}
\endhead
\bottomrule\noalign{}
\endlastfoot
FR-001--002, D1 & Antes de atender una operación protegida, el servidor
deberá comprobar sesión vigente, capacidad y ámbito aplicables. Si no se
satisfacen, rechazará la operación sin modificar el recurso. & Todos;
sesión ausente, válida y fuera de ámbito. \\
FR-005, D1 & Un administrador autorizado podrá habilitar o desactivar
una identidad y conservará su autoría histórica; la desactivación
impedirá acceso posterior. & Administrador; desactivar y repetir
solicitud con sesión previa. \\
FR-043, D4 & El administrador funcional podrá registrar carreras con
código normalizado único, nombre, unidad académica y vigencia válida; se
rechazará el código duplicado y se conservará el historial de cambios. &
Administrador; alta válida, duplicado y cambio histórico. \\
FR-007/045, D1/D4 & Una versión de modelo admitirá nodos tipados con
distinta profundidad; rechazará ciclos y pertenencias inválidas antes de
publicarse. & Administrador de modelos; indicador directo/anidado y
ciclo. \\
FR-008--009, D1 & La publicación validará estructura, fuentes,
metodología y revisión independiente; una versión publicada no podrá
sobrescribirse y su corrección originará una sucesora. & Autor/revisor;
publicación incompleta y sucesión. \\
FR-011/048, D1/D4 & El coordinador podrá preparar un proceso para
carrera, período y modelo; la activación exigirá carrera/modelo
aplicables y autorización registrada y verificada. &
Coordinador/autoridad; sin autorización permanece PREPARATION. \\
FR-013/050, D1/D4 & La planificación asignará participantes o equipos
con función, indicador y vigencia explícitos; rechazará asignaciones
incompatibles sin excepción válida. & Coordinador; vigencia, conflicto y
ámbito. \\
FR-016--018/051, D1/D4 & El responsable podrá preparar evidencia de su
ámbito y declarar coberturas exactas; remitir exigirá metadatos
completos y contenido compatible con cada fuente cubierta. &
Responsable; fuente ajena, contenido incompleto y caso válido. \\
FR-019--020, D1 & Una remisión conservará versión, autor y fecha; no
admitirá edición posterior. Una corrección creará una sucesora sin
eliminar la predecesora. & Responsable; remitir, intentar edición y
corregir. \\
FR-021--022, D1 & El revisor autorizado registrará resultados por
cobertura de una versión remitida y se derivará el estado global de esos
resultados, conservando actor, fecha y fundamento aplicable. & Revisor;
aceptar una cobertura y observar otra. \\
FR-023/055, D1/D4 & Aceptar documentalmente una evidencia no creará una
valoración académica ni marcará el indicador como cumplido. &
Revisor/evaluador; comparar registros académicos antes/después. \\
FR-053, D4 y D2 & La recopilación será el número de IDs distintos de
fuentes obligatorias aplicables con cobertura remitida vigente dividido
entre las fuentes obligatorias aplicables, por cien; borradores y
fuentes opcionales no contarán. & Coordinador; 3 de 4 = 75 \%. \\
FR-054, D4 y D2 & La revisión documental contará fuentes obligatorias
aplicables cuya cardinalidad se satisfaga con coberturas aceptadas
vigentes; no duplicará una fuente por tener varios archivos. &
Coordinador; 2 de 4 = 50 \% y cardinalidad mayor que uno. \\
FR-030, D1 & Desde un agregado autorizado, el usuario podrá acceder a
indicadores, decisiones y versiones que explican el conteo, conservando
las restricciones de ámbito. & Consultor/coordinador; cada destino abre
el sustento correspondiente. \\
FR-031/064, D1/D4 & El sistema distinguirá un corte sin cambio de estado
del cierre documental sujeto a política; ninguno exigirá ejecutar
valoración, informes o mejora posteriores. & Coordinador; corte con
pendientes y MVP sin P12--P15. \\
FR-039, D1 & Los errores expondrán una estructura estable con código,
descripción, instancia y correlación, sin secretos ni información de
ámbitos ajenos. & Todos; entrada inválida y error de negocio. \\
\end{longtable}
}

Observaciones: FR-030 mantiene implementación parcial en el tablero
localizado. La consulta/registro de decisiones existe, pero no demuestra
toda la política institucional. Los requisitos compuestos deben
desglosarse al diseñar pruebas; aprobar un caso no verifica
automáticamente todas sus cláusulas.

\section{Requisitos no funcionales y mecanismo de
medición}\label{requisitos-no-funcionales-y-mecanismo-de-mediciuxf3n}

Los umbrales numéricos provienen del repositorio. Las condiciones
experimentales pendientes se identifican, sin inventar resultados.

{\def\LTcaptype{none} % do not increment counter
\begin{longtable}[]{@{}
  >{\raggedright\arraybackslash}p{(\linewidth - 4\tabcolsep) * \real{0.3333}}
  >{\raggedright\arraybackslash}p{(\linewidth - 4\tabcolsep) * \real{0.3333}}
  >{\raggedright\arraybackslash}p{(\linewidth - 4\tabcolsep) * \real{0.3333}}@{}}
\toprule\noalign{}
\begin{minipage}[b]{\linewidth}\raggedright
Categoría / ID
\end{minipage} & \begin{minipage}[b]{\linewidth}\raggedright
Exigencia y forma de verificar
\end{minipage} & \begin{minipage}[b]{\linewidth}\raggedright
Estado y límite
\end{minipage} \\
\midrule\noalign{}
\endhead
\bottomrule\noalign{}
\endlastfoot
Seguridad, RNF-001 & Almacenar contraseñas mediante hash adaptativo con
sal; inspeccionar almacenamiento y verificar autenticación sin exponer
secretos. & BCrypt localizado; no es auditoría integral de seguridad. \\
Seguridad, RNF-002 & Expirar sesión tras 30 minutos de inactividad;
cerrar sesión o desactivar cuenta invalidará acceso posterior. Verificar
con reloj controlado y solicitudes sucesivas. & Mecanismos y pruebas
existentes. \\
Seguridad, RNF-003 & Limitar autenticaciones al alcanzar cinco fallos en
15 minutos. Comparar error de usuario inexistente y contraseña
incorrecta. & Confirmar si son consecutivos o acumulados: texto y SQL
difieren. \\
Rendimiento, RNF-006 & p95 de operaciones interactivas del MVP ≤ 2 s,
excluyendo archivos e informes. Definir inicio/fin visibles, entorno,
red y operaciones antes de medir. & Medición histórica de consultas
seleccionadas: no acredita todos los recorridos visibles. \\
Rendimiento/capacidad, RNF-007--008 & p95 de cargas hasta 25 MiB ≤ 15 s;
al menos 100 usuarios, 30 sesiones, 500 indicadores, 10 000 versiones y
100 GiB sin incumplir RNF-006. & Falta campaña de cargas y volumen 100
GiB acreditado. \\
Disponibilidad, RNF-009 & Un fallo de exportación/notificación futura no
impedirá autenticación ni operaciones del MVP. Ejecutar recorridos con
ese componente fallando y comparar estado. & No se exige un porcentaje
de uptime no definido; ensayo de fallos pendiente. \\
Recuperación, RNF-010 & Una interrupción entre metadatos y archivos
dejará operación completa o recuperable, sin archivo ausente presentado
como disponible. Interrumpir, reconciliar y comprobar integridad. &
Recuperación localizada; RPO/RTO institucional por definir. \\
Usabilidad, SC-003/007 & Recorrido D1--D6 en ≤ 45 min; al menos 19 de 20
participantes completan la tarea al primer intento según protocolo. &
Los datos deben ser sintéticos; los participantes humanos
representativos. Protocolo/resultados pendientes. \\
Accesibilidad, RNF-014 & Recorridos declarados con teclado, foco
visible, etiquetas y contraste WCAG 2.2 AA. Verificar criterios
aplicables automáticamente y con revisión humana. & Recorridos Edge no
certifican todo el nivel AA. \\
Mantenibilidad, RNF-017--018 & Mantener reglas separadas de
controladores, persistencia e interfaz; verificar dependencias y pruebas
reproducibles de estados, permisos y cálculo. & ArchUnit y suites
disponibles; dependencias SQL requieren revisión adicional. \\
Interoperabilidad, FR-035 y contratos & Validar formatos versionados,
entradas y errores al intercambiar información; distinguir contrato
operativo y futuro. & No existe RNF independiente específico;
integración externa aún futura. \\
\end{longtable}
}

\textbf{Candidato no aprobado de interoperabilidad (ACD-INT-01):} ``Los
clientes de las operaciones HTTP seleccionadas deberán poder interpretar
las representaciones JSON y los errores documentados, conservando
nombres, tipos y semántica del contrato versionado''. Verificación
propuesta: cliente HTTP independiente y pruebas de esquema contra rutas
operativas. Se deriva del diseño REST; el equipo debe validarlo antes de
integrarlo al catálogo RNF oficial. No obliga a habilitar CORS ni a
implementar importación externa futura.

\section{Calidad y requisitos
corregidos}\label{calidad-y-requisitos-corregidos}

{\def\LTcaptype{none} % do not increment counter
\begin{longtable}[]{@{}
  >{\raggedright\arraybackslash}p{(\linewidth - 4\tabcolsep) * \real{0.3333}}
  >{\raggedright\arraybackslash}p{(\linewidth - 4\tabcolsep) * \real{0.3333}}
  >{\raggedright\arraybackslash}p{(\linewidth - 4\tabcolsep) * \real{0.3333}}@{}}
\toprule\noalign{}
\begin{minipage}[b]{\linewidth}\raggedright
ID
\end{minipage} & \begin{minipage}[b]{\linewidth}\raggedright
Redacción o problema de partida
\end{minipage} & \begin{minipage}[b]{\linewidth}\raggedright
Formulación revisada y justificación
\end{minipage} \\
\midrule\noalign{}
\endhead
\bottomrule\noalign{}
\endlastfoot
FR-053 & ``Fuentes aplicables'' puede incluir opcionales y duplicar
coberturas. & Contar IDs distintos de fuentes obligatorias aplicables
cubiertas por versión remitida vigente. Hace explícitos conjunto,
denominador y vigencia; coincide con D2 y C9. \\
FR-030 & ``Permitir navegar'' no define origen, destino ni comprobación.
& Desde cada agregado autorizado se accederá a sus indicadores,
decisiones y versiones exactas, manteniendo ámbito. Permite comprobar
cada enlace y detectar una interfaz incompleta. \\
RNF-006 & ``Resultado visible'' no equivale a duración del endpoint. &
Mantener p95 ≤ 2 s y definir medición desde acción del usuario hasta
presentación del resultado, con conjunto de operaciones y entorno
registrado. Evita usar una medición API como prueba del recorrido
completo. \\
RNF-003 & El enunciado dice fallos consecutivos; el SQL cuenta fallos
dentro de ventana. & Pendiente de decisión: conservar consecutividad y
ajustar comportamiento, o aprobar ventana acumulada de cinco fallos/15
min. No se elige una política silenciosamente. \\
SC-007 & Se menciona ``participantes sintéticos''. & Participantes
representativos utilizando datos sintéticos. La usabilidad requiere
personas; la protección de datos se mantiene. \\
\end{longtable}
}

Comprobación de los cinco criterios de la guía:

{\def\LTcaptype{none} % do not increment counter
\begin{longtable}[]{@{}
  >{\raggedright\arraybackslash}p{(\linewidth - 10\tabcolsep) * \real{0.1667}}
  >{\raggedright\arraybackslash}p{(\linewidth - 10\tabcolsep) * \real{0.1667}}
  >{\raggedright\arraybackslash}p{(\linewidth - 10\tabcolsep) * \real{0.1667}}
  >{\raggedright\arraybackslash}p{(\linewidth - 10\tabcolsep) * \real{0.1667}}
  >{\raggedright\arraybackslash}p{(\linewidth - 10\tabcolsep) * \real{0.1667}}
  >{\raggedright\arraybackslash}p{(\linewidth - 10\tabcolsep) * \real{0.1667}}@{}}
\toprule\noalign{}
\begin{minipage}[b]{\linewidth}\raggedright
Requisito revisado
\end{minipage} & \begin{minipage}[b]{\linewidth}\raggedright
Claro
\end{minipage} & \begin{minipage}[b]{\linewidth}\raggedright
Específico
\end{minipage} & \begin{minipage}[b]{\linewidth}\raggedright
Pertinente
\end{minipage} & \begin{minipage}[b]{\linewidth}\raggedright
Verificable
\end{minipage} & \begin{minipage}[b]{\linewidth}\raggedright
Consistente con alcance
\end{minipage} \\
\midrule\noalign{}
\endhead
\bottomrule\noalign{}
\endlastfoot
FR-053 & Define qué se cuenta. & Distingue obligatorias y vigentes. &
Seguimiento documental. & Ejemplo 3/4 y cero fuentes. & P11, MVP. \\
FR-030 & Identifica navegación y sustentos. & Define destinos y ámbito.
& Trazabilidad. & Apertura y permisos de cada destino. & MVP;
implementación parcial. \\
RNF-006 & Resultado visible definido. & p95 y exclusiones existentes. &
Respuesta del sistema. & Requiere protocolo de entorno. & MVP; no amplía
entregas. \\
RNF-003 & Aún requiere elegir política. & Ventana/umbral conocidos. &
Control de acceso. & Dos secuencias de intentos. & Acceso; no aprobado
el cambio. \\
SC-007 & Persona y datos diferenciados. & 19 de 20 participantes. &
Usabilidad. & Registro de tareas y ayuda. & Recorridos documentales. \\
\end{longtable}
}

\section{Casos de uso
representativos}\label{casos-de-uso-representativos}

Los identificadores P proceden de D3; se detallan cuatro escenarios
suficientes para el propósito de ACD2.

\subsection{P02 --- Registrar una
carrera}\label{p02-registrar-una-carrera}

\begin{itemize}
\tightlist
\item
  Actor principal: administrador funcional autorizado.
\item
  Objetivo: disponer de una carrera identificable para posteriores
  procesos.
\item
  Precondiciones: sesión vigente, capacidad de administración y unidad
  académica activa.
\item
  Flujo principal: seleccionar unidad; introducir código, nombre y
  vigencia; enviar; normalizar código; validar unicidad/unidad/fechas;
  guardar historial y auditoría; responder con carrera, Location y ETag.
\item
  Alternativa: código repetido produce 409 y no crea otra carrera.
\item
  Excepción: solicitud malformada o campos inválidos produce 400; unidad
  no disponible produce 422 según C2.
\item
  Postcondición de éxito: carrera activa con revisión inicial 1, unidad
  y autoría conservadas.
\item
  Postcondición de rechazo: no se crea la carrera solicitada.
\item
  Requisitos: FR-043, FR-002, FR-006; RNF-011--012.
\end{itemize}

\subsection{P08 --- Remitir una versión de
evidencia}\label{p08-remitir-una-versiuxf3n-de-evidencia}

\begin{itemize}
\tightlist
\item
  Actor principal: responsable documental.
\item
  Objetivo: poner una versión completa a disposición del revisor.
\item
  Precondiciones: proceso habilitado, asignación vigente, borrador
  completo, contenido permitido y coberturas exactas.
\item
  Flujo principal: consultar versión/revisión; comprobar metadatos,
  fuentes y contenido; remitir con revisión vigente; reautorizar y
  validar integridad; conservar versión inmutable y autor/fecha;
  incorporarla a revisión.
\item
  Alternativa: metadatos o fuentes incompletos conservan el borrador y
  bloquean remisión.
\item
  Excepción: revisión obsoleta produce 412; no debe sobrescribirse la
  versión ajena.
\item
  Postcondición: versión SUBMITTED inmutable o borrador conservado sin
  remisión parcial.
\item
  Requisitos: FR-018--019, FR-051; RNF-010--012.
\end{itemize}

\subsection{P09 --- Revisar coberturas
documentales}\label{p09-revisar-coberturas-documentales}

\begin{itemize}
\tightlist
\item
  Actor principal: revisor documental autorizado.
\item
  Objetivo: determinar aceptabilidad documental de cada cobertura.
\item
  Precondiciones: versión remitida vigente, capacidad de revisión y
  separación de funciones aplicable.
\item
  Flujo principal: consultar contenido; revisar cada fuente; registrar
  ACCEPTED/OBSERVED/REJECTED y fundamento; derivar global; conservar
  decisión; actualizar proyección documental.
\item
  Alternativa: aceptar una cobertura y observar otra produce global
  OBSERVED; únicamente la cobertura aceptada puede contribuir conforme a
  cardinalidad.
\item
  Excepción: autorrevisión sin excepción válida o acceso no autorizado
  se rechazan; la política de excepción exige validación adicional.
\item
  Postcondición: decisión conservada y sin valoración académica creada
  automáticamente.
\item
  Requisitos: FR-021--024, FR-054--055.
\end{itemize}

\subsection{P11 --- Consultar seguimiento
documental}\label{p11-consultar-seguimiento-documental}

\begin{itemize}
\tightlist
\item
  Actor principal: coordinador o consultor autorizado.
\item
  Objetivo: identificar pendientes sin confundir documentación con
  cumplimiento académico.
\item
  Precondiciones: proceso existente y ámbito de consulta vigente.
\item
  Flujo principal: filtrar fuentes autorizadas; excluir opcionales/no
  aplicables del denominador; contar remisiones y aceptaciones vigentes;
  calcular porcentajes; mostrar sustento y política.
\item
  Alternativa: cuatro fuentes, tres remitidas y dos aceptadas producen
  75 \% y 50 \%.
\item
  Excepción/caso límite: sin configuración devuelve porcentaje no
  calculable; todas excluidas autorizadamente tampoco equivalen a 100
  \%.
\item
  Postcondición: consulta sin cambios de estado y agregados
  reproducibles.
\item
  Requisitos: FR-029--030, FR-053--054.
\end{itemize}

\section{Historias de usuario y criterios de
aceptación}\label{historias-de-usuario-y-criterios-de-aceptaciuxf3n}

Son selecciones/refinamientos académicos de historias existentes; no se
atribuye su aprobación a la sesión.

\subsection{HU-ACD-01 --- Catálogo de carreras
(005/US1)}\label{hu-acd-01-catuxe1logo-de-carreras-005us1}

Como administrador funcional, quiero registrar carreras con su unidad y
vigencia, para preparar procesos sobre un catálogo identificable y
conservar su historial.

\begin{itemize}
\tightlist
\item
  AC-01: \textbf{Dado que} la unidad está activa y el código no existe,
  \textbf{cuando} registro campos válidos, \textbf{entonces} se crea una
  carrera con revisión 1 y se devuelve 201 y su Location.
\item
  AC-02: \textbf{Dado que} el código normalizado ya existe,
  \textbf{cuando} intento registrar otro igual, \textbf{entonces} recibo
  409 y no se crea un duplicado.
\item
  AC-03: \textbf{Dado que} mi identidad no tiene capacidad de
  administración, \textbf{cuando} intento crear, \textbf{entonces} la
  operación se rechaza y no modifica el catálogo.
\end{itemize}

\subsection{HU-ACD-02 --- Remisión trazable (004/US3 y
005/US4)}\label{hu-acd-02-remisiuxf3n-trazable-004us3-y-005us4}

Como responsable documental, quiero remitir una versión completa
vinculada a fuentes exactas, para que el revisor conozca qué
documentación se aporta y bajo qué versión.

\begin{itemize}
\tightlist
\item
  AC-04: \textbf{Dado que} el borrador satisface metadatos, contenido y
  cobertura, \textbf{cuando} lo remito con revisión vigente,
  \textbf{entonces} queda SUBMITTED con autor y fecha y no admite
  edición ordinaria.
\item
  AC-05: \textbf{Dado que} falta una cobertura requerida,
  \textbf{cuando} intento remitir, \textbf{entonces} no se crea una
  remisión parcial y se conserva el borrador.
\item
  AC-06: \textbf{Dado que} otro cambio avanzó la revisión,
  \textbf{cuando} envío un If-Match obsoleto, \textbf{entonces} recibo
  412 y debo recargar antes de decidir.
\end{itemize}

\subsection{HU-ACD-03 --- Revisión por cobertura (005/US5, parte del
revisor)}\label{hu-acd-03-revisiuxf3n-por-cobertura-005us5-parte-del-revisor}

Como revisor documental, quiero decidir por cobertura de una versión
remitida, para comunicar su aceptabilidad sin atribuir un resultado
académico.

\begin{itemize}
\tightlist
\item
  AC-07: \textbf{Dado que} hay dos coberturas, \textbf{cuando} acepto
  una y observo otra con fundamento, \textbf{entonces} se conservan
  ambos resultados y el global se deriva como OBSERVED.
\item
  AC-08: \textbf{Dado que} no existe valoración académica,
  \textbf{cuando} acepto documentación, \textbf{entonces} no se crea una
  valoración ni se marca cumplimiento académico.
\end{itemize}

\subsection{HU-ACD-04 --- Seguimiento reproducible
(005/US6)}\label{hu-acd-04-seguimiento-reproducible-005us6}

Como coordinador, quiero consultar conteos y porcentajes separados por
fuentes, para intervenir sobre pendientes documentales y explicar el
sustento de cada agregado.

\begin{itemize}
\tightlist
\item
  AC-09: \textbf{Dado que} hay cuatro fuentes obligatorias aplicables,
  tres con remisión vigente y dos con cardinalidad aceptada,
  \textbf{cuando} consulto seguimiento, \textbf{entonces} obtengo 75 \%
  de recopilación y 50 \% de revisión.
\item
  AC-10: \textbf{Dado que} una fuente opcional tiene varios archivos,
  \textbf{cuando} consulto, \textbf{entonces} no aumenta el denominador
  obligatorio ni se cuenta varias veces.
\item
  AC-11: \textbf{Dado que} no existen fuentes obligatorias configuradas,
  \textbf{cuando} consulto, \textbf{entonces} no se divide entre cero ni
  se informa 100 \%.
\end{itemize}

\section{Trazabilidad de la
selección}\label{trazabilidad-de-la-selecciuxf3n}

{\def\LTcaptype{none} % do not increment counter
\begin{longtable}[]{@{}
  >{\raggedright\arraybackslash}p{(\linewidth - 6\tabcolsep) * \real{0.2500}}
  >{\raggedright\arraybackslash}p{(\linewidth - 6\tabcolsep) * \real{0.2500}}
  >{\raggedright\arraybackslash}p{(\linewidth - 6\tabcolsep) * \real{0.2500}}
  >{\raggedright\arraybackslash}p{(\linewidth - 6\tabcolsep) * \real{0.2500}}@{}}
\toprule\noalign{}
\begin{minipage}[b]{\linewidth}\raggedright
Necesidad
\end{minipage} & \begin{minipage}[b]{\linewidth}\raggedright
Requisitos
\end{minipage} & \begin{minipage}[b]{\linewidth}\raggedright
Historia / caso
\end{minipage} & \begin{minipage}[b]{\linewidth}\raggedright
Implementación y prueba existente
\end{minipage} \\
\midrule\noalign{}
\endhead
\bottomrule\noalign{}
\endlastfoot
Identificar la carrera y preservar historial & FR-043 & HU-ACD-01 / P02
& C1--C3; \path|tests/acceptance/career-catalog.spec.ts|. \\
Conservar información remitida exacta & FR-018--020/051 & HU-ACD-02 /
P08/P10 &
\path|backend/src/main/java/ec/edu/unl/sec/evidence/application/SubmitEvidenceService.java|;
\path|tests/acceptance/document-review-correction.spec.ts|. \\
Separar dictamen documental de académico & FR-021--023/055 & HU-ACD-03 /
P09 &
\path|backend/src/main/java/ec/edu/unl/sec/documentreview/application/DocumentReviewService.java|;
\path|tests/acceptance/complete-workflows.spec.ts|, AC-P09-01. \\
Explicar avance por fuentes & FR-030/053/054 & HU-ACD-04 / P11 & C9;
suite consolidada AC-P11-01. FR-030 sigue parcial en interfaz. \\
\end{longtable}
}

Las pruebas se identifican como existentes, no como ejecutadas
nuevamente. Las tres pruebas vacías \path|evidence-v1.spec.ts|,
\path|evidence-v1-v2.spec.ts| y \path|evidence-audit.spec.ts| no se
utilizan como prueba de cumplimiento.

\chapter{Contraste técnico: HTTP, JSON y revisión
final}\label{contraste-tuxe9cnico-http-json-y-revisiuxf3n-final}

\section{Funcionamiento básico y
arquitectura}\label{funcionamiento-buxe1sico-y-arquitectura}

HTTP es un protocolo de aplicación de solicitud/respuesta. El método
expresa la operación, la URL identifica el recurso, las cabeceras
aportan metadatos y el contenido transporta una representación cuando
corresponde. El servidor responde con estado, cabeceras y contenido
opcional. La semántica del protocolo se mantiene aunque la aplicación
conserve una sesión mediante cookies. Fuente:
\href{https://www.rfc-editor.org/rfc/rfc9110.html}{RFC 9110}.

En este proyecto, Angular utiliza rutas relativas \path|/api/v1|
mediante C12; Spring Boot recibe las solicitudes, autoriza y coordina
módulos, persistencia PostgreSQL y archivos. El backend es monolítico
modular; no se afirma que existan microservicios ni un broker activo.

Una solicitud comprende línea de petición, cabeceras, línea en blanco y
cuerpo opcional. En los ejemplos se usa notación HTTP/1.1 legible; no se
afirma que una captura real negociara esa versión. Una URL completa
distingue esquema, host, puerto, ruta y consulta. Por ejemplo,
\path|http://localhost:8080/api/v1/careers?page=0&size=20| identifica el
servicio local, catálogo y parámetros de paginación; el fragmento de URL
no se envía al servidor.

\section{Comparación de GET, POST, PUT, PATCH y
DELETE}\label{comparaciuxf3n-de-get-post-put-patch-y-delete}

{\def\LTcaptype{none} % do not increment counter
\begin{longtable}[]{@{}
  >{\raggedright\arraybackslash}p{(\linewidth - 4\tabcolsep) * \real{0.3333}}
  >{\raggedright\arraybackslash}p{(\linewidth - 4\tabcolsep) * \real{0.3333}}
  >{\raggedright\arraybackslash}p{(\linewidth - 4\tabcolsep) * \real{0.3333}}@{}}
\toprule\noalign{}
\begin{minipage}[b]{\linewidth}\raggedright
Método
\end{minipage} & \begin{minipage}[b]{\linewidth}\raggedright
Semántica general
\end{minipage} & \begin{minipage}[b]{\linewidth}\raggedright
Ejemplo real del proyecto
\end{minipage} \\
\midrule\noalign{}
\endhead
\bottomrule\noalign{}
\endlastfoot
GET & Recuperar una representación; seguro e idempotente. No se solicita
una mutación de negocio. & \path|/api/v1/careers/{id}| devuelve carrera
y ETag. \\
POST & Procesar contenido conforme al recurso; puede crear o ejecutar
una acción. No es idempotente por definición. & \path|/api/v1/careers|
crea; \path|/api/v1/sessions| inicia sesión. \\
PUT & Crear/reemplazar el estado representado del recurso objetivo;
semántica idempotente. & \path|/api/v1/model-versions/{id}| reemplaza un
borrador con el DTO legado. \\
PATCH & Aplicar modificaciones a un recurso; no es necesariamente
idempotente. & \path|/api/v1/users/{userId}| recibe
\path|{active:false}|. \\
DELETE & Solicitar retirar la asociación del recurso objetivo; semántica
idempotente. & Retirar una referencia externa del borrador; logout en
\path|/api/v1/sessions|. \\
\end{longtable}
}

Fuentes:
\href{https://developer.mozilla.org/en-US/docs/Web/HTTP/Reference/Methods}{MDN:
métodos} y \href{https://www.rfc-editor.org/rfc/rfc5789.html}{RFC 5789:
PATCH}.

Idempotencia significa que repetir una operación tiene el mismo efecto
pretendido, no que siempre devuelva el mismo código. Un segundo DELETE
puede ser rechazado o encontrar el recurso ausente. Un segundo PUT con
revisión vieja puede recibir 412 sin contradecir esa semántica. PATCH no
significa enviar cualquier subconjunto de campos: manda el contrato
concreto. Por ejemplo, el PATCH de carreras exige todos los campos de
CareerInput; no se debe omitir \path|academicUnitId| suponiendo un JSON
Merge Patch genérico.

\section{Códigos de estado
solicitados}\label{cuxf3digos-de-estado-solicitados}

{\def\LTcaptype{none} % do not increment counter
\begin{longtable}[]{@{}
  >{\raggedright\arraybackslash}p{(\linewidth - 4\tabcolsep) * \real{0.3333}}
  >{\raggedright\arraybackslash}p{(\linewidth - 4\tabcolsep) * \real{0.3333}}
  >{\raggedright\arraybackslash}p{(\linewidth - 4\tabcolsep) * \real{0.3333}}@{}}
\toprule\noalign{}
\begin{minipage}[b]{\linewidth}\raggedright
Estado
\end{minipage} & \begin{minipage}[b]{\linewidth}\raggedright
Interpretación general
\end{minipage} & \begin{minipage}[b]{\linewidth}\raggedright
Ejemplo respaldado por el código
\end{minipage} \\
\midrule\noalign{}
\endhead
\bottomrule\noalign{}
\endlastfoot
200 OK & Solicitud atendida correctamente. & GET de carrera o PUT del
borrador. \\
201 Created & Se crearon uno o más recursos. & POST de carrera con
Location. \\
204 No Content & Operación completada sin contenido de respuesta. &
DELETE de referencia o cierre de sesión. \\
400 Bad Request & Solicitud no procesable por error del cliente. &
JSON/UUID inválido o campos no válidos. \\
401 Unauthorized & Faltan credenciales válidas para el recurso. & Sesión
ausente, revocada o expirada. \\
403 Forbidden & El servidor rechaza la operación. & Sesión válida sin
capacidad o token CSRF inválido. \\
404 Not Found & No se encuentra o no se revela el recurso. & Carrera
desconocida o no visible. \\
409 Conflict & Conflicto con el estado actual. & Código de carrera
duplicado. \\
422 Unprocessable Content & Contenido comprendido, instrucciones no
procesables. & Unidad académica no disponible, según C2. \\
500 Internal Server Error & Condición inesperada del servidor. & C10
devuelve INTERNAL\_ERROR sin detalle técnico al cliente. \\
\end{longtable}
}

Fuente de semántica general:
\href{https://developer.mozilla.org/en-US/docs/Web/HTTP/Reference/Status}{MDN:
estados HTTP}. La asociación con rutas concretas procede del
repositorio.

También son relevantes 412 (revisión obsoleta), 428 (falta If-Match),
413 (archivo demasiado grande) y 429 (limitación de autenticación). Un
error de negocio no debe interpretarse automáticamente como 500; el
cuerpo explica el código de dominio. Un fallo de CORS del navegador
tampoco equivale necesariamente a un 500 del backend.

\section{Cabeceras y autenticación real del
proyecto}\label{cabeceras-y-autenticaciuxf3n-real-del-proyecto}

{\def\LTcaptype{none} % do not increment counter
\begin{longtable}[]{@{}
  >{\raggedright\arraybackslash}p{(\linewidth - 4\tabcolsep) * \real{0.3333}}
  >{\raggedright\arraybackslash}p{(\linewidth - 4\tabcolsep) * \real{0.3333}}
  >{\raggedright\arraybackslash}p{(\linewidth - 4\tabcolsep) * \real{0.3333}}@{}}
\toprule\noalign{}
\begin{minipage}[b]{\linewidth}\raggedright
Cabecera
\end{minipage} & \begin{minipage}[b]{\linewidth}\raggedright
Propósito
\end{minipage} & \begin{minipage}[b]{\linewidth}\raggedright
Aplicación o límite
\end{minipage} \\
\midrule\noalign{}
\endhead
\bottomrule\noalign{}
\endlastfoot
Content-Type & Describe el formato del contenido transportado. &
\path|application/json|; para errores \path|application/problem+json|;
los archivos se cargan con multipart. \\
Accept & Indica formatos de respuesta admitidos por el cliente. & Puede
indicar JSON y Problem Details; no convierte por sí sola el cuerpo
enviado. \\
Authorization & Transporta credenciales según un esquema de
autenticación. & \path|Bearer <token>| es una forma común, pero no es el
mecanismo implementado de SIGEAC. \\
Location & Referencia un recurso creado o destino, según la respuesta. &
En 201 de carrera: \path|/api/v1/careers/{id}|. No significa una
redirección automática. \\
ETag & Validador opaco de una representación. & C11 utiliza revisión
fuerte entre comillas, por ejemplo \path|"1"|. No es credencial de
autorización. \\
If-Match & Condiciona la modificación al validador recibido. & Evita
sobrescritura silenciosa; si está obsoleto, 412. \\
Cookie / Set-Cookie & Envío y establecimiento de estado de sesión. &
Cookie SEC\_SESSION; no enviar su valor real en el informe. \\
X-CSRF-TOKEN & Token propio de protección de mutaciones autenticadas. &
C5 valida su hash; una cookie válida no sustituye el token. \\
\end{longtable}
}

Referencias:
\href{https://developer.mozilla.org/en-US/docs/Web/HTTP/Reference/Headers/Content-Type}{Content-Type},
\href{https://developer.mozilla.org/en-US/docs/Web/HTTP/Reference/Headers/Accept}{Accept},
\href{https://developer.mozilla.org/en-US/docs/Web/HTTP/Reference/Headers/Authorization}{Authorization},
\href{https://developer.mozilla.org/en-US/docs/Web/HTTP/Reference/Headers/Location}{Location}
y
\href{https://developer.mozilla.org/en-US/docs/Web/HTTP/Reference/Headers/ETag}{ETag}.

El login de C4 devuelve \path|csrfToken| y establece SEC\_SESSION. No
devuelve JWT ni habilita \path|Authorization: Bearer| por el ejemplo de
la guía. Se conserva ese concepto en el análisis general y se usa Cookie
en los ejemplos del proyecto. En C14 Spring desactiva CSRF integrado,
pero C5 implementa una comprobación propia; no se concluye ausencia de
protección a partir de una sola línea.

\section{Ejercicio resuelto: POST de carrera y respuesta
JSON}\label{ejercicio-resuelto-post-de-carrera-y-respuesta-json}

Ejemplo sintético compatible con C1--C3. Los UUID, fecha y tokens no
representan personas ni solicitudes reales. Se presupone una unidad
activa con el ID indicado y un administrador autorizado.

\begin{Verbatim}[breaklines=true,breakanywhere=true,fontsize=\footnotesize]
POST /api/v1/careers HTTP/1.1
Host: localhost:8080
Content-Type: application/json
Accept: application/json, application/problem+json
Cookie: SEC_SESSION=<sesion-sintetica>
X-CSRF-TOKEN: <csrf-sintetico>

{
  "code": "CARRERA-DEMO",
  "name": "Carrera de demostración",
  "academicUnitId": "11111111-1111-4111-8111-111111111111",
  "validFrom": "2026-01-01",
  "validUntil": null
}
\end{Verbatim}

\begin{Verbatim}[breaklines=true,breakanywhere=true,fontsize=\footnotesize]
HTTP/1.1 201 Created
Content-Type: application/json
Location: /api/v1/careers/22222222-2222-4222-8222-222222222222
ETag: "1"

{
  "id": "22222222-2222-4222-8222-222222222222",
  "data": {
    "code": "CARRERA-DEMO",
    "name": "Carrera de demostración",
    "academicUnitId": "11111111-1111-4111-8111-111111111111",
    "validFrom": "2026-01-01",
    "validUntil": null
  },
  "status": "ACTIVE",
  "revision": 1,
  "createdBy": "33333333-3333-4333-8333-333333333333",
  "createdAt": "2026-01-01T12:00:00Z",
  "classificationPending": false
}
\end{Verbatim}

{\def\LTcaptype{none} % do not increment counter
\begin{longtable}[]{@{}
  >{\raggedright\arraybackslash}p{(\linewidth - 2\tabcolsep) * \real{0.5000}}
  >{\raggedright\arraybackslash}p{(\linewidth - 2\tabcolsep) * \real{0.5000}}@{}}
\toprule\noalign{}
\begin{minipage}[b]{\linewidth}\raggedright
Elemento
\end{minipage} & \begin{minipage}[b]{\linewidth}\raggedright
Interpretación
\end{minipage} \\
\midrule\noalign{}
\endhead
\bottomrule\noalign{}
\endlastfoot
Método & POST solicita crear una carrera mediante el catálogo. \\
Recurso & \path|/api/v1/careers| es la colección; la carrera creada
tiene URI propia. \\
URL y Host & La ruta y \path|Host| permiten identificar el servicio y
recurso. \\
Cabeceras de petición & Content-Type describe JSON; Accept admite
éxito/error; Cookie identifica sesión; CSRF protege la mutación. \\
Cuerpo & Código, nombre, unidad y vigencia; el servidor asigna ID,
autor, estado y revisión. \\
Estado & 201 es apropiado porque se creó un recurso, no solo se consultó
uno. \\
Location & Identifica la nueva carrera; C1 implementa el GET de esa
ruta. \\
ETag & Revisión inicial \path|"1"| para futuras modificaciones
condicionadas. \\
Cuerpo de respuesta & Representación del recurso; \path|data| contiene
los atributos de negocio. \\
\end{longtable}
}

Si se repite el código, C2 lanza 409 DUPLICATE\_CODE; no se promete que
repetir POST cree o devuelva siempre el mismo objeto. Si la unidad no
está activa, se produce 422 UNIT\_UNAVAILABLE. El cuerpo debe respetar
la distinción de fechas ISO, cadenas y \path|null|.

\section{GET y PATCH: consulta y
concurrencia}\label{get-y-patch-consulta-y-concurrencia}

\begin{Verbatim}[breaklines=true,breakanywhere=true,fontsize=\footnotesize]
GET /api/v1/careers/22222222-2222-4222-8222-222222222222 HTTP/1.1
Host: localhost:8080
Accept: application/json
Cookie: SEC_SESSION=<sesion-sintetica>
\end{Verbatim}

La respuesta exitosa es 200 con la representación Career y ETag. No se
envía cuerpo para este GET.

\begin{Verbatim}[breaklines=true,breakanywhere=true,fontsize=\footnotesize]
PATCH /api/v1/careers/22222222-2222-4222-8222-222222222222 HTTP/1.1
Host: localhost:8080
Content-Type: application/json
Cookie: SEC_SESSION=<sesion-sintetica>
X-CSRF-TOKEN: <csrf-sintetico>
If-Match: "1"

{
  "code": "CARRERA-DEMO",
  "name": "Carrera de demostración actualizada",
  "academicUnitId": "11111111-1111-4111-8111-111111111111",
  "validFrom": "2026-01-01",
  "validUntil": null
}
\end{Verbatim}

Con revisión 1 vigente, la respuesta es 200, revisión 2 y ETag
\path|"2"|. Si otro cliente ya guardó la revisión 2, el PATCH anterior
recibe 412 y no sobrescribe sus cambios. La respuesta de error tiene
forma Problem Details; este ejemplo es reconstruido desde C2/C16:

\begin{Verbatim}[breaklines=true,breakanywhere=true,fontsize=\footnotesize]
HTTP/1.1 412 Precondition Failed
Content-Type: application/problem+json

{
  "type": "urn:sec-unl:problem:stale-revision",
  "title": "Precondition Failed",
  "status": 412,
  "detail": "El recurso cambió; recargue antes de continuar.",
  "instance": "/api/v1/careers/22222222-2222-4222-8222-222222222222",
  "code": "STALE_REVISION",
  "correlationId": "44444444-4444-4444-8444-444444444444",
  "details": {}
}
\end{Verbatim}

No confundir 412 con 409: el primero señala que no se cumplió la
condición de versión; el segundo puede señalar código duplicado o estado
incompatible.

\section{PUT y DELETE implementados}\label{put-y-delete-implementados}

PUT operativo: \path|/api/v1/model-versions/{id}|. El DTO legado tiene
\path|source|, \path|criteria|, \path|scale| y \path|formula|; el DTO
configurado de \path|/configured| tiene \path|source|, \path|nodes| y
\path|configuration|. Debe enviarse la representación admitida completa,
sesión, CSRF e If-Match. Éxito: 200 y ETag. Un modelo publicado requiere
sucesión; no se presupone que el PUT cree un recurso inexistente porque
la semántica general lo permita.

DELETE operativo: retirar una referencia externa de un borrador
autorizado:

\begin{Verbatim}[breaklines=true,breakanywhere=true,fontsize=\footnotesize]
DELETE /api/v1/evidence-versions/55555555-5555-4555-8555-555555555555/external-references/66666666-6666-4666-8666-666666666666 HTTP/1.1
Host: localhost:8080
Cookie: SEC_SESSION=<sesion-sintetica>
X-CSRF-TOKEN: <csrf-sintetico>
If-Match: "3"
\end{Verbatim}

\begin{Verbatim}[breaklines=true,breakanywhere=true,fontsize=\footnotesize]
HTTP/1.1 204 No Content
ETag: "4"
\end{Verbatim}

No hay cuerpo JSON tras 204. Se retira la referencia del borrador, no
necesariamente el contenido externo al que apunta. C8 llama al servicio
que exige versión editable; no es un permiso para borrar versiones
remitidas.

\section{JSON: interpretación de datos y
errores}\label{json-interpretaciuxf3n-de-datos-y-errores}

JSON representa objetos, arreglos, cadenas, números, booleanos y
\path|null|; los nombres de objeto y cadenas usan comillas dobles. No
incorpora comentarios ni comas finales. Fechas e identificadores se
representan como cadenas con convenciones del contrato. Fuente:
\href{https://www.rfc-editor.org/rfc/rfc8259.html}{RFC 8259}.

En el alta de carrera, \path|revision| es número,
\path|classificationPending| booleano y \path|validUntil| puede ser
\path|null|; \path|"null"| sería una cadena distinta. Un JSON
sintácticamente válido todavía puede infringir una regla de negocio: por
ejemplo, una unidad inexistente produce 422 en C2. Un cuerpo malformado
se interpreta como 400 en C10. Los campos de autoría/revisión los
controla el servidor y no se añaden al cuerpo CareerInput.

\section{CORS y orígenes distintos}\label{cors-y-oruxedgenes-distintos}

Un origen combina esquema, host y puerto. \path|http://localhost:4200| y
\path|http://localhost:8080| son orígenes distintos. CORS permite al
servidor declarar qué orígenes pueden leer respuestas desde el
navegador; no sustituye autenticación ni autorización. Fuente:
\href{https://fetch.spec.whatwg.org/\#http-cors-protocol}{Fetch
Standard: CORS}.

En C12 el cliente usa \path|/api/v1|; C13 configura proxy \path|/api|
hacia \path|http://localhost:8080|. Si el navegador accede a Angular en
4200 y solicita esa ruta relativa, ve el mismo origen; el proxy realiza
el salto al backend. No se localizaron políticas CORS explícitas en las
fuentes revisadas y no se afirma que el acceso directo entre orígenes
esté habilitado.

Si se despliegan frontend y API en orígenes distintos, deben definirse
orígenes permitidos, métodos y cabeceras; con cookies, credenciales y
política de cookies compatibles. Con credenciales no se usa
\path|Access-Control-Allow-Origin: *|. Para leer ETag o Location desde
JavaScript pueden requerirse \path|Access-Control-Expose-Headers|. JSON
y cabeceras como If-Match/CSRF pueden provocar preflight OPTIONS.
Referencia:
\href{https://developer.mozilla.org/en-US/docs/Web/HTTP/Guides/CORS}{MDN:
CORS}.

Ejercicio hipotético de preflight, \textbf{no implementado ni ejecutado
como política del proyecto}:

\begin{Verbatim}[breaklines=true,breakanywhere=true,fontsize=\footnotesize]
OPTIONS /api/v1/careers HTTP/1.1
Host: localhost:8080
Origin: http://localhost:4200
Access-Control-Request-Method: POST
Access-Control-Request-Headers: content-type,x-csrf-token
\end{Verbatim}

Para permitir ese escenario, una configuración futura tendría que
responder con origen explícito, método/cabeceras permitidos y
credenciales cuando corresponda. Un navegador puede impedir leer la
respuesta si faltan permisos CORS, aunque la API funcione en Postman. No
se modifica la configuración como parte de ACD2.

\section{Socialización y
retroalimentación}\label{socializaciuxf3n-y-retroalimentaciuxf3n}

No se ha aportado evidencia de la sesión. El análisis técnico siguiente
sirve como insumo de discusión, no como testimonio de compañeros o
docente.

{\def\LTcaptype{none} % do not increment counter
\begin{longtable}[]{@{}
  >{\raggedright\arraybackslash}p{(\linewidth - 4\tabcolsep) * \real{0.3333}}
  >{\raggedright\arraybackslash}p{(\linewidth - 4\tabcolsep) * \real{0.3333}}
  >{\raggedright\arraybackslash}p{(\linewidth - 4\tabcolsep) * \real{0.3333}}@{}}
\toprule\noalign{}
\begin{minipage}[b]{\linewidth}\raggedright
Observación técnica preparada
\end{minipage} & \begin{minipage}[b]{\linewidth}\raggedright
Mejora incorporada a este desarrollo
\end{minipage} & \begin{minipage}[b]{\linewidth}\raggedright
Validación grupal
\end{minipage} \\
\midrule\noalign{}
\endhead
\bottomrule\noalign{}
\endlastfoot
Un archivo no equivale a una fuente aceptada ni a cumplimiento
académico. & FR-053/054 distinguen fuentes, remisiones y aceptación;
AC-08/09 comprueban separación. & {[}PENDIENTE: confirmar discusión y
acuerdo{]} \\
``Rápido'' no ofrece un criterio de aprobación. & RNF-006 conserva p95 ≤
2 s y exige protocolo del resultado visible. & {[}PENDIENTE{]} \\
El ejemplo Bearer de la guía no coincide con la aplicación. & Ejercicios
usan SEC\_SESSION y CSRF, sin inventar JWT. & {[}PENDIENTE{]} \\
Una cabecera ETag sin If-Match no evita por sí sola sobrescrituras. &
Ejercicio de concurrencia interpreta 412 y recarga. & {[}PENDIENTE{]} \\
\end{longtable}
}

Registro para la sesión real:

{\def\LTcaptype{none} % do not increment counter
\begin{longtable}[]{@{}
  >{\raggedright\arraybackslash}p{(\linewidth - 8\tabcolsep) * \real{0.2000}}
  >{\raggedright\arraybackslash}p{(\linewidth - 8\tabcolsep) * \real{0.2000}}
  >{\raggedright\arraybackslash}p{(\linewidth - 8\tabcolsep) * \real{0.2000}}
  >{\raggedright\arraybackslash}p{(\linewidth - 8\tabcolsep) * \real{0.2000}}
  >{\raggedright\arraybackslash}p{(\linewidth - 8\tabcolsep) * \real{0.2000}}@{}}
\toprule\noalign{}
\begin{minipage}[b]{\linewidth}\raggedright
Persona que observa
\end{minipage} & \begin{minipage}[b]{\linewidth}\raggedright
Observación recibida y fecha
\end{minipage} & \begin{minipage}[b]{\linewidth}\raggedright
Requisito/ejercicio afectado
\end{minipage} & \begin{minipage}[b]{\linewidth}\raggedright
Cambio realizado
\end{minipage} & \begin{minipage}[b]{\linewidth}\raggedright
Evidencia/acuerdo
\end{minipage} \\
\midrule\noalign{}
\endhead
\bottomrule\noalign{}
\endlastfoot
{[}PENDIENTE{]} & {[}PENDIENTE: observación real{]} & {[}PENDIENTE{]} &
{[}PENDIENTE{]} & {[}PENDIENTE{]} \\
\end{longtable}
}

No debe marcarse esta tabla como completada copiando una observación
hipotética. La evidencia individual complementaria contiene dos
requisitos corregidos y un ejercicio desarrollado, pero también exige el
comentario real de retroalimentación de cada estudiante.

\section{Comprobación de la guía y
conclusiones}\label{comprobaciuxf3n-de-la-guuxeda-y-conclusiones}

{\def\LTcaptype{none} % do not increment counter
\begin{longtable}[]{@{}
  >{\raggedright\arraybackslash}p{(\linewidth - 4\tabcolsep) * \real{0.3333}}
  >{\raggedright\arraybackslash}p{(\linewidth - 4\tabcolsep) * \real{0.3333}}
  >{\raggedright\arraybackslash}p{(\linewidth - 4\tabcolsep) * \real{0.3333}}@{}}
\toprule\noalign{}
\begin{minipage}[b]{\linewidth}\raggedright
Criterio de la actividad
\end{minipage} & \begin{minipage}[b]{\linewidth}\raggedright
Evidencia preparada
\end{minipage} & \begin{minipage}[b]{\linewidth}\raggedright
Pendiente real
\end{minipage} \\
\midrule\noalign{}
\endhead
\bottomrule\noalign{}
\endlastfoot
Problema, actores, alcance y clarificaciones & Secciones 1.2--1.5,
fuentes D1--D4. & Diagnóstico institucional y respuestas verificadas. \\
Requisitos funcionales & Sección 2.1 y verificaciones. & Aprobación del
equipo y decisiones abiertas. \\
Siete categorías de RNF & Sección 2.2, umbrales y límites. & Protocolos
y resultados que faltan. \\
Redacción clara y verificable & Sección 2.3 y justificaciones. &
Registrar aceptación de correcciones. \\
Casos de uso e historias & Secciones 2.4--2.6; AC-01--11. &
Socialización y ajustes acordados. \\
HTTP y JSON & Secciones 3.1--3.8; solicitudes/respuestas resueltas. & Si
se pide ejecución, capturas reales sanitizadas. \\
CORS & Sección 3.9 y proxy existente. & Decidir política solo si el
despliegue lo requiere. \\
Evidencia individual & \path|EVIDENCIA_INDIVIDUAL.md|. & Nombre, aporte
real, retroalimentación y revisión personal. \\
Participación y retroalimentación & Registro preparado en 3.10. & No
acreditada todavía. \\
\end{longtable}
}

El aporte del análisis consiste en vincular necesidades documentadas con
condiciones y resultados comprobables: conservar versiones, autorizar
por ámbito y contar fuentes distintas. Los ejemplos HTTP muestran que el
mismo recurso puede consultarse, crearse o modificarse con métodos y
precondiciones diferentes. ETag/If-Match permite explicar concurrencia
entre clientes; Cookie y CSRF reflejan la autenticación real del
sistema. El informe no acredita participación, ejecución ni aprobación
institucional inexistentes.

\cleardoublepage
\nocite{proyecto,rfc9110,rfc5789,rfc8259,fetch,mdnmethods,mdnstatus,mdncors,mdnetag,mdnauthorization,mdncontenttype,mdnaccept,mdnlocation}
\bibliographystyle{IEEEtran}
\bibliography{referencias}
\appendix
\chapter{Evidencia individual para
personalizar}\label{evidencia-individual-para-personalizar}

{\def\LTcaptype{none} % do not increment counter
\begin{longtable}[]{@{}
  >{\raggedright\arraybackslash}p{(\linewidth - 2\tabcolsep) * \real{0.5000}}
  >{\raggedright\arraybackslash}p{(\linewidth - 2\tabcolsep) * \real{0.5000}}@{}}
\toprule\noalign{}
\begin{minipage}[b]{\linewidth}\raggedright
Dato
\end{minipage} & \begin{minipage}[b]{\linewidth}\raggedright
Valor
\end{minipage} \\
\midrule\noalign{}
\endhead
\bottomrule\noalign{}
\endlastfoot
Estudiante & {[}PENDIENTE: nombre completo{]} \\
Equipo y ciclo/paralelo & {[}PENDIENTE{]} \\
Fecha de participación & {[}PENDIENTE: fecha real{]} \\
Actividad & ACD2 --- Semana 2, Desarrollo de Software Distribuido \\
Proyecto & SIGEAC / SEC-UNL \\
Estado & Desarrollo técnico propuesto; pendiente de revisión y
atribución personal \\
\end{longtable}
}

Este contenido desarrolla los ejercicios, pero no acredita por sí solo
participación individual. Cada estudiante debe revisarlo, explicar su
razonamiento y registrar su retroalimentación real. No presentar el
mismo texto como aporte personal independiente de todos los integrantes
sin identificar la colaboración y el uso de IA.

\section{Dos requisitos revisados y
justificación}\label{dos-requisitos-revisados-y-justificaciuxf3n}

\subsection{FR-053 --- Cálculo de recopilación
documental}\label{fr-053-cuxe1lculo-de-recopilaciuxf3n-documental}

\textbf{Fuente:} \path|specs/005-complete-evaluation-workflows/spec.md|,
complementada por \path|docs/requirements/domain-rules.md|, sección
Cálculo documental sin ambigüedad.

\textbf{Enunciado de partida:} ``El avance de recopilación DEBE
calcularse como fuentes aplicables cubiertas por evidencia remitida
dividido para fuentes requeridas aplicables, multiplicado por cien''.

\textbf{Redacción corregida propuesta:} El avance de recopilación deberá
calcularse como el número de IDs distintos de fuentes obligatorias
aplicables cubiertas por al menos una versión remitida vigente, dividido
entre el número de IDs distintos de fuentes obligatorias aplicables,
multiplicado por cien. Se excluirán borradores y fuentes opcionales. Si
el denominador es cero, se mostrará porcentaje no calculable y el estado
correspondiente de configuración o no aplicabilidad; no se mostrará 100
\%.

\textbf{Justificación:} ``Fuentes aplicables'' no precisa por sí sola
obligatoriedad, duplicación ni vigencia. La corrección conserva la
intención y utiliza reglas existentes, haciendo posible diseñar casos de
prueba sin depender del número de archivos.

\textbf{Criterio:} Dado que existen cuatro fuentes obligatorias
aplicables y tres tienen cobertura remitida vigente, cuando consulto
recopilación, entonces obtengo 75 \%. Adjuntar dos archivos a una misma
fuente no debe aumentar ese conteo.

\textbf{Evidencia técnica:}
\path|backend/src/main/java/ec/edu/unl/sec/reporting/domain/DocumentaryProgressCalculator.java|;
escenario AC-P11-01 de
\path|tests/acceptance/complete-workflows.spec.ts|. No se ejecutó la
prueba durante la preparación documental.

\subsection{FR-030 --- Navegación hacia los
sustentos}\label{fr-030-navegaciuxf3n-hacia-los-sustentos}

\textbf{Fuente:} \path|docs/requirements/spec.md|.

\textbf{Enunciado de partida:} ``El tablero DEBE permitir navegar desde
conteos agregados hasta indicadores, decisiones y evidencias autorizadas
que los sustentan''.

\textbf{Redacción corregida propuesta:} Desde cada conteo agregado
presentado en el tablero, el usuario autorizado deberá poder acceder a
los indicadores, decisiones documentales y versiones de evidencia
exactas que forman el conteo. Cada consulta de destino mantendrá el
ámbito del usuario y no expondrá registros ajenos. La interfaz deberá
conservar el contexto del proceso consultado.

\textbf{Justificación:} Se especifican origen, destinos y restricción
observable. Así se distingue una navegación funcional de mostrar un UUID
sin enlace. La conservación de contexto se propone para revisión del
equipo y no se presenta como una decisión previamente aprobada.

\textbf{Criterio:} Dado que un conteo incluye una cobertura aceptada de
mi ámbito, cuando accedo a su sustento, entonces veo el indicador,
decisión y versión exacta relacionados. Dado que carezco de autorización
para un recurso, cuando intento abrirlo, entonces no recibo su contenido
ni se amplía mi ámbito.

\textbf{Evidencia y límite:}
\path|frontend/src/app/features/reporting/pages/documentary-dashboard.page.ts|
filtra por fuente, pero muestra UUID de versiones como texto; la
navegación completa sigue parcial. Formular el requisito no demuestra
que esté implementado.

\textbf{Aporte personal a confirmar:} {[}PENDIENTE: explicar con
palabras propias qué parte se revisó, qué decisión se discutió y cuál
fue el cambio aceptado por el grupo{]}.

\section{Ejercicio de análisis HTTP
resuelto}\label{ejercicio-de-anuxe1lisis-http-resuelto}

Ejemplo reconstruido del contrato de
\path|backend/src/main/java/ec/edu/unl/sec/selfassessment/api/CareerController.java|
con datos sintéticos. Se presupone sesión administrativa vigente y
unidad activa. No es una captura de ejecución real.

\begin{Verbatim}[breaklines=true,breakanywhere=true,fontsize=\footnotesize]
POST /api/v1/careers HTTP/1.1
Host: localhost:8080
Content-Type: application/json
Accept: application/json
Cookie: SEC_SESSION=<sesion-sintetica>
X-CSRF-TOKEN: <csrf-sintetico>

{
  "code": "CARRERA-DEMO",
  "name": "Carrera de demostración",
  "academicUnitId": "11111111-1111-4111-8111-111111111111",
  "validFrom": "2026-01-01",
  "validUntil": null
}
\end{Verbatim}

\begin{Verbatim}[breaklines=true,breakanywhere=true,fontsize=\footnotesize]
HTTP/1.1 201 Created
Content-Type: application/json
Location: /api/v1/careers/22222222-2222-4222-8222-222222222222
ETag: "1"

{
  "id": "22222222-2222-4222-8222-222222222222",
  "data": {
    "code": "CARRERA-DEMO",
    "name": "Carrera de demostración",
    "academicUnitId": "11111111-1111-4111-8111-111111111111",
    "validFrom": "2026-01-01",
    "validUntil": null
  },
  "status": "ACTIVE",
  "revision": 1,
  "createdBy": "33333333-3333-4333-8333-333333333333",
  "createdAt": "2026-01-01T12:00:00Z",
  "classificationPending": false
}
\end{Verbatim}

{\def\LTcaptype{none} % do not increment counter
\begin{longtable}[]{@{}
  >{\raggedright\arraybackslash}p{(\linewidth - 2\tabcolsep) * \real{0.5000}}
  >{\raggedright\arraybackslash}p{(\linewidth - 2\tabcolsep) * \real{0.5000}}@{}}
\toprule\noalign{}
\begin{minipage}[b]{\linewidth}\raggedright
Elemento solicitado
\end{minipage} & \begin{minipage}[b]{\linewidth}\raggedright
Interpretación
\end{minipage} \\
\midrule\noalign{}
\endhead
\bottomrule\noalign{}
\endlastfoot
Método & POST: se solicita crear una carrera. \\
Recurso & Colección \path|/api/v1/careers|; el servidor asigna una URI a
la carrera nueva. \\
Host & \path|localhost:8080| identifica el servicio usado en el
ejemplo. \\
Content-Type & El cuerpo de la petición es JSON; la respuesta también
declara JSON. \\
Accept & El cliente solicita una representación JSON de respuesta. \\
Cookie & Lleva la sesión opaca; el servidor comprueba su vigencia y
permisos. \\
X-CSRF-TOKEN & Protege la mutación autenticada contra solicitudes
cruzadas no autorizadas. \\
Cuerpo de petición & Código, nombre, ID de unidad y vigencia. No envía
autor ni revisión. \\
Estado & 201 Created confirma creación; no es simplemente una consulta
200. \\
Location & Identifica la carrera recién creada y permite construir una
consulta posterior. \\
ETag & \path|"1"| es la revisión inicial; se usa con If-Match en cambios
posteriores. \\
Cuerpo de respuesta & Representación Career, incluyendo atributos de
negocio en \path|data| y metadatos del servidor. \\
Significado completo & Se creó una carrera activa, con unidad y vigencia
válidas, bajo una identidad autorizada. \\
\end{longtable}
}

\textbf{Comparación de rechazos:} sin sesión corresponde 401; sin
capacidad o con CSRF inválido, 403; código duplicado, 409; unidad no
disponible, 422 según el servicio. No todos los JSON erróneos producen
422: errores sintácticos o de validación de campos pueden producir 400.

\textbf{Pregunta de concurrencia resuelta:} si leo ETag \path|"1"| y
otro cliente modifica antes, mi PATCH con If-Match \path|"1"| recibe
412; debo recargar. ETag no autentica mi identidad y no autoriza mi
acción.

\textbf{Pregunta CORS resuelta:} el cliente Angular usa rutas relativas
y un proxy de desarrollo. Si llama directamente desde un origen distinto
al backend, el navegador requiere una política CORS compatible; Postman
no verifica ese comportamiento del navegador. No se presume CORS
habilitado.

\textbf{Authorization de la guía:} su ejemplo usa Bearer; este sistema
utiliza SEC\_SESSION y CSRF. Se interpreta Authorization como cabecera
general sin afirmar soporte JWT en el proyecto.

Fuentes conceptuales:
\href{https://www.rfc-editor.org/rfc/rfc9110.html}{RFC 9110},
\href{https://developer.mozilla.org/en-US/docs/Web/HTTP/Reference/Headers/ETag}{MDN:
ETag} y \href{https://fetch.spec.whatwg.org/\#http-cors-protocol}{Fetch:
CORS}. Las rutas, DTO y códigos aplicados proceden del repositorio.

\section{Comentario de retroalimentación
real}\label{comentario-de-retroalimentaciuxf3n-real}

No se proporcionó una observación recibida en clase. Completar el
siguiente registro después de la socialización; no atribuir la revisión
de IA al docente o a un compañero.

{\def\LTcaptype{none} % do not increment counter
\begin{longtable}[]{@{}
  >{\raggedright\arraybackslash}p{(\linewidth - 2\tabcolsep) * \real{0.5000}}
  >{\raggedright\arraybackslash}p{(\linewidth - 2\tabcolsep) * \real{0.5000}}@{}}
\toprule\noalign{}
\begin{minipage}[b]{\linewidth}\raggedright
Campo
\end{minipage} & \begin{minipage}[b]{\linewidth}\raggedright
Registro
\end{minipage} \\
\midrule\noalign{}
\endhead
\bottomrule\noalign{}
\endlastfoot
Persona que realizó la observación & {[}PENDIENTE{]} \\
Fecha y contexto & {[}PENDIENTE{]} \\
Observación recibida & {[}PENDIENTE: texto fiel o resumen
identificado{]} \\
Requisito o ejercicio afectado & {[}PENDIENTE{]} \\
Cambio utilizado para mejorar el trabajo & {[}PENDIENTE{]} \\
Justificación de cómo resolvió el problema & {[}PENDIENTE{]} \\
Evidencia de participación & {[}PENDIENTE: acta, revisión, entrega o
registro real{]} \\
\end{longtable}
}

Estructura de redacción para completar: ``Durante la sesión,
{[}persona{]} observó que {[}observación real{]}. A partir de ello
revisé {[}ID o ejercicio{]} y cambié {[}antes{]} por {[}después{]},
porque {[}justificación verificable{]}''.

\section{Comprobación individual}\label{comprobaciuxf3n-individual}

\begin{itemize}
\tightlist
\item
  Dos requisitos revisados con justificación: desarrollados en 1.1--1.2;
  revisión y adopción personal pendientes.
\item
  Solicitud/respuesta interpretadas con método, recurso, cabeceras,
  cuerpo y estado: desarrolladas en 2.
\item
  Retroalimentación real y uso explicado: pendiente en 3.
\item
  Nombre, fecha y contribución individual: pendientes.
\item
  Las solicitudes son sintéticas y las pruebas no se presentan como
  ejecutadas.
\end{itemize}

\end{document}
